\chapter{Introdu{\c c}{\~a}o {\`a} Avalia{\c c}{\~a}o Formativa e {\`a} IA}
\label{cap:introducao}

\section{O Prop{\'o}sito deste Manual e do Reposit{\'o}rio}

O reposit{\'o}rio \texttt{skills-ia-educacao} constitui um acervo documental e tecnol{\'o}gico concebido para preencher a lacuna entre as diretrizes normativas (como o Referencial do MEC \cite{brasil2026referencial} e as resolu{\c c}{\~o}es universit{\'a}rias, por exemplo, da UNIFEI \cite{unifei2025graduacao}) e a pr{\'a}tica cotidiana do corpo docente. Trata-se de produto do \textbf{Projeto de Ensino AvalIA -- Kit de Avalia{\c c}{\~a}o Formativa com Intelig{\^e}ncia Artificial para Professores do Ensino Superior}, desenvolvido com o apoio da Pr{\'o}-Reitoria de Gradua{\c c}{\~a}o da UNIFEI: sua finalidade {\'e} instrumentar o corpo docente com recursos prontos de avalia{\c c}{\~a}o formativa mediada por IA, materializando as diretrizes institucionais em ferramentas de uso imediato.

Em vez de exigir que cada professor se torne um especialista em engenharia de \emph{prompts} complexos, o reposit{\'o}rio disponibiliza \textbf{62 skills tem{\'a}ticas}, reutiliz{\'a}veis em CLIs compat{\'i}veis com o padr{\~a}o Agent Skills, e \textbf{6 subagentes} configurados especificamente para operar no \textbf{OpenCode} (\url{https://opencode.ai}).

\begin{conceitochave}
\textbf{OpenCode, skill e subagente em tr{\^e}s frases}: pense no \textbf{OpenCode} como a oficina (o programa que executa as tarefas); na \textbf{skill} como a ficha de instru{\c c}{\~o}es de uma tarefa (ex.: \texttt{/ia-educacao-rubrica} ensina a montar uma rubrica); e no \textbf{subagente} como o especialista que segue v{\'a}rias fichas para entregar um trabalho completo (ex.: \texttt{@construtor-de-avaliacoes} monta prova com rubrica e gabarito). Voc{\^e} n{\~a}o precisa decorar nada: basta descrever o que quer ou digitar a barra (\texttt{/}) para skills e o arroba (\texttt{@}) para subagentes.
\end{conceitochave}

\begin{objetivos}
\begin{itemize}
  \item Compreender por que a avalia{\c c}{\~a}o formativa se torna central quando a IA generativa transforma as tarefas acad{\^e}micas;
  \item Conhecer os cinco n{\'i}veis da Escala AIAS e diferenciar seu car{\'a}ter n{\~a}o hier{\'a}rquico de outras taxonomias;
  \item Visualizar o percurso completo do manual e saber a que cap{\'i}tulo recorrer em cada necessidade docente.
\end{itemize}
\end{objetivos}

\section{Rota R{\'a}pida para Come{\c c}ar}

Se esta {\'e} a sua primeira experi{\^e}ncia com IA generativa, n{\~a}o {\'e} necess{\'a}rio dominar todo o manual antes de realizar a primeira atividade. O percurso abaixo leva cerca de 60 minutos e termina com um material rascunhado, revisado e com a pol{\'i}tica de IA declarada. Siga nesta ordem:

\begin{enumerate}
  \item Leia este cap{\'i}tulo e o Cap{\'i}tulo~\ref{cap:fundamentos} para compreender a finalidade pedag{\'o}gica da IA (cerca de 15 minutos; leitura conceitual, sem computador);
  \item Instale o OpenCode, configure um provedor e instale as skills e os subagentes seguindo o Cap{\'i}tulo~\ref{cap:instalacao} (cerca de 20 minutos; use apenas o \emph{caminho recomendado} do seu sistema e ignore os m{\'e}todos alternativos por enquanto);
  \item Execute um caso simples do Cap{\'i}tulo~\ref{cap:casos-de-uso}, substituindo a disciplina e o tema pelos seus (cerca de 15 minutos; o Caso~2 {\'e} uma boa primeira escolha para quem avalia, o Caso~1 para quem planeja);
  \item Revise o resultado com o checklist do pr{\'o}prio Cap{\'i}tulo~\ref{cap:casos-de-uso} antes de entregar qualquer material aos estudantes (cerca de 5 minutos; nada vai aos estudantes sem passar por esta revis{\~a}o);
  \item Use o Cap{\'i}tulo~\ref{cap:explicar-estudantes} para declarar a pol{\'i}tica de uso de IA da turma (cerca de 5 minutos; copie o texto pronto de declara{\c c}{\~a}o). Os Cap{\'i}tulos~\ref{cap:skills-opencode} e~\ref{cap:agentes-skills} s{\~a}o refer{\^e}ncia de consulta: podem ficar para depois, quando voc{\^e} quiser entender ou manter o acervo.
\end{enumerate}

\subsection{Vocabul{\'a}rio M{\'i}nimo}

\begin{itemize}
  \item \textbf{IAGen}: intelig{\^e}ncia artificial generativa, capaz de produzir texto, c{\'o}digo, imagens ou outros artefatos a partir de instru{\c c}{\~o}es;
  \item \textbf{LLM}: modelo de linguagem de grande escala, usado para prever e gerar sequ{\^e}ncias de texto;
  \item \textbf{Prompt}: instru{\c c}{\~a}o, contexto e crit{\'e}rios fornecidos ao modelo;
  \item \textbf{Skill}: arquivo de instru{\c c}{\~o}es especializadas carregado sob demanda;
  \item \textbf{Subagente}: especialista configurado para coordenar uma tarefa complexa e, conforme suas permiss{\~o}es, usar v{\'a}rias skills;
  \item \textbf{AIAS}: escala que explicita como a IA pode participar de uma atividade avaliativa;
  \item \textbf{STHEM}: sigla para as grandes {\'a}reas Science, Technology, Humanities, Engineering e Mathematics (Ci{\^e}ncias, Tecnologia, Humanidades, Engenharias e Matem{\'a}tica);
  \item \textbf{DUA}: abordagem que reduz barreiras oferecendo m{\'u}ltiplos meios de engajamento, representa{\c c}{\~a}o e express{\~a}o;
  \item \textbf{CoVe}: procedimento de verifica{\c c}{\~a}o independente de afirma{\c c}{\~o}es produzidas por IA;
  \item \textbf{CLI e TUI}: \emph{CLI} {\'e} o terminal (tela escura de comandos); \emph{TUI} {\'e} a interface interativa que abre ao digitar \texttt{opencode} no terminal;
  \item \textbf{Provedor e modelo}: o \emph{provedor} {\'e} a empresa ou servi{\c c}o que oferece a IA; o \emph{modelo} {\'e} o ``c{\'e}rebro'' escolhido na hora (ver Cap{\'i}tulo~\ref{cap:modelos-gratuitos});
  \item \textbf{Headless}: modo de usar o OpenCode com uma {\'u}nica linha de comando, sem abrir a interface (ex.: \texttt{opencode run "..."});
  \item \textbf{Frontmatter}: pequeno cabe{\c c}alho t{\'e}cnico no in{\'i}cio de cada skill; voc{\^e} n{\~a}o precisa edit{\'a}-lo para usar as skills;
  \item \textbf{MCP}: padr{\~a}o aberto que liga o OpenCode a ferramentas externas; citado aqui s{\'o} para registro de proced{\^e}ncia;
  \item \textbf{ConcepTest e blueprint}: \emph{ConcepTest} {\'e} uma quest{\~a}o conceitual curta para vota{\c c}{\~a}o em sala; \emph{blueprint} {\'e} o mapa da prova (quantas quest{\~o}es por t{\'o}pico e n{\'i}vel).
\end{itemize}

Ao longo dos pr{\'o}ximos cap{\'i}tulos, este manual apresentar{\'a} em detalhes:
\begin{itemize}
  \item Os fundamentos pedag{\'o}gicos de feedback, taxonomia cognitiva e rubricas anal{\'i}ticas (Cap{\'i}tulo~\ref{cap:fundamentos});
  \item A arquitetura das skills e seu mecanismo de execu{\c c}{\~a}o com OpenCode (Cap{\'i}tulo~\ref{cap:skills-opencode});
  \item O passo a passo completo de instala{\c c}{\~a}o do OpenCode e sincroniza{\c c}{\~a}o do acervo (Cap{\'i}tulo~\ref{cap:instalacao});
  \item O fluxo de trabalho na interface gr{\'a}fica e desktop (Cap{\'i}tulo~\ref{cap:desktop});
  \item A escolha de modelos gratuitos no OpenCode (Cap{\'i}tulo~\ref{cap:modelos-gratuitos});
  \item Princ{\'i}pios de {\'e}tica, limita{\c c}{\~o}es, LGPD e transpar{\^e}ncia (Cap{\'i}tulo~\ref{cap:etica-transparencia});
  \item Estrat{\'e}gias para explicar a IA e a Escala AIAS aos estudantes (Cap{\'i}tulo~\ref{cap:explicar-estudantes});
  \item A especifica{\c c}{\~a}o dos seis subagentes e o cat{\'a}logo das 62 skills (Cap{\'i}tulo~\ref{cap:agentes-skills});
  \item Casos de uso iniciais e roteiros pr{\'a}ticos para sala de aula (Cap{\'i}tulo~\ref{cap:casos-de-uso}).
\end{itemize}

\section{O Desafio Pedag{\'o}gico na Era da Intelig{\^e}ncia Artificial}

A r{\'a}pida difus{\~a}o de modelos de linguagem de grande escala (\emph{Large Language Models} -- LLMs) e de ferramentas baseadas em Intelig{\^e}ncia Artificial Generativa (IAGen) vem transformando a educa{\c c}{\~a}o b{\'a}sica e superior. Tarefas cl{\'a}ssicas de avalia{\c c}{\~a}o, como reda{\c c}{\~a}o de ensaios curtos, resolu{\c c}{\~a}o de quest{\~o}es conceituais elementares e gera{\c c}{\~a}o de trechos introdut{\'o}rios de c{\'o}digo-fonte, podem ser realizadas com facilidade por assistentes computacionais \cite{unesco2021guidance,mollick2024instructors}.

Frente a esse cen{\'a}rio, rea{\c c}{\~o}es puramente proibitivas ou tentativas ing{\^e}nuas de detec{\c c}{\~a}o automatizada de conte{\'u}do gerado por IA demonstraram-se ineficazes, eticamente arriscadas e pedagogicamente est{\'e}reis. Como alerta o Minist{\'e}rio da Educa{\c c}{\~a}o em seu documento norteador, o \emph{Referencial para Desenvolvimento e Uso Respons{\'a}veis de Intelig{\^e}ncia Artificial na Educa{\c c}{\~a}o} \cite{brasil2026referencial}, a intelig{\^e}ncia artificial deve ser compreendida n{\~a}o como uma amea{\c c}a a ser banida, mas como um catalisador de revis{\~a}o das pr{\'a}ticas docentes, exigindo centralidade humana, transpar{\^e}ncia {\'e}tica e intencionalidade pedag{\'o}gica.

A superada {\^e}nfase em produtos finais padronizados precisa ceder espa{\c c}o para a valoriza{\c c}{\~a}o dos processos de pensamento, da investiga{\c c}{\~a}o cr{\'i}tica, da metacogni{\c c}{\~a}o e da avalia{\c c}{\~a}o formativa cont{\'i}nua.

\section{Da Avalia{\c c}{\~a}o Somativa {\`a} Avalia{\c c}{\~a}o Formativa}

Historicamente, as institui{\c c}{\~o}es de ensino superior priorizaram a \textbf{avalia{\c c}{\~a}o somativa} -- instrumentos pontuais aplicados ao t{\'e}rmino de uma unidade ou semestre (provas tradicionais, exames finais e testes de m{\'u}ltipla escolha sem devolutiva) cujo objetivo prec{\'i}puo {\'e} certificar, classificar ou atribuir uma nota ao estudante. 

Embora a certifica{\c c}{\~a}o cumpra papel institucional regulat{\'o}rio, a literatura cient{\'i}fica internacional h{\'a} d{\'e}cadas demonstra que a avalia{\c c}{\~a}o somativa isolada possui impacto limitado sobre a qualidade efetiva da aprendizagem \cite{black1998assessment}. Em contrapartida, a \textbf{avalia{\c c}{\~a}o formativa} concebe a avalia{\c c}{\~a}o como parte intr{\'i}nseca do ato de ensinar e aprender. 

\begin{conceitochave}
\textbf{Avalia{\c c}{\~a}o Formativa}: Qualquer processo avaliativo planejado e implementado ao longo do percurso de aprendizagem com a finalidade de coletar evid{\^e}ncias sobre a compreens{\~a}o dos estudantes, fornecendo informa{\c c}{\~o}es que permitam a professores e alunos ajustarem o ensino e as estrat{\'e}gias de estudo em tempo real \cite{black1998assessment}.
\end{conceitochave}

Segundo a formula{\c c}{\~a}o seminal de \citeonline{sadler1989formative}, a avalia{\c c}{\~a}o formativa eficaz exige que o estudante compreenda tr{\^e}s elementos fundamentais:
\begin{enumerate}
  \item \textbf{O padr{\~a}o ou objetivo de refer{\^e}ncia}: clareza sobre o que constitui um desempenho de excel{\^e}ncia;
  \item \textbf{O estado atual de aprendizagem}: reconhecimento cr{\'i}tico de onde o pr{\'o}prio trabalho se situa no momento;
  \item \textbf{As a{\c c}{\~o}es de fechamento do hiato (\emph{closing the gap})}: estrat{\'e}gias claras e acion{\'a}veis para mover o desempenho atual em dire{\c c}{\~a}o ao padr{\~a}o desejado.
\end{enumerate}

Em ambientes educacionais permeados por IA, essa concep{\c c}{\~a}o ganha urg{\^e}ncia. Se uma m{\'a}quina pode redigir a resposta final, o valor pedag{\'o}gico desloca-se para a capacidade do aluno em interrogar o problema, planejar abordagens, validar solu{\c c}{\~o}es, corrigir inconsist{\^e}ncias e justificar decis{\~o}es tomadas.

\section{A Escala AIAS: Transpar{\^e}ncia e Intencionalidade Pedag{\'o}gica}

Para que a integra{\c c}{\~a}o da intelig{\^e}ncia artificial n{\~a}o ocorra em um v{\'a}cuo normativo ou sob regras amb{\'i}guas, a literatura contempor{\^a}nea desenvolveu a \textbf{AIAS} (\emph{Artificial Intelligence Assessment Scale}), proposta originalmente por \citeonline{perkins2024aias} e refinada por \citeonline{perkins2025reimagining}.

A AIAS {\'e} um referencial conceitual que classifica o n{\'i}vel de uso de IA admitido ou exigido em uma determinada atividade avaliativa. Ela n{\~a}o se destina a vigiar o aluno, mas sim a explicitar o contrato pedag{\'o}gico entre docente e discente. A Tabela~\ref{tab:aias-escala} apresenta os cinco n{\'i}veis can{\^o}nicos adotados estritamente em todo o ecossistema deste reposit{\'o}rio.

\begin{table}[htbp]
  \centering
  \small
  \caption{Os 5 N{\'i}veis Can{\^o}nicos da Escala AIAS}
  \label{tab:aias-escala}
  \begin{tabularx}{\linewidth}{clp{4cm}X}
    \toprule
    \textbf{N{\'i}vel} & \textbf{Nome Can{\^o}nico} & \textbf{IA Permitida} & \textbf{Produto Final Esperado} \\
    \midrule
    \textbf{1} & Sem IA & Nenhuma interven{\c c}{\~a}o computacional inteligente & Elaborado 100\% pelo estudante sem qualquer aux{\'i}lio de IA. \\
    \addlinespace
    \textbf{2} & Planejamento Assistido por IA & Idea{\c c}{\~a}o, \emph{brainstorming}, esbo{\c c}os e estrutura{\c c}{\~a}o & Do estudante; a IA atua apenas na fase inicial preparat{\'o}ria. \\
    \addlinespace
    \textbf{3} & Colabora{\c c}{\~a}o com IA & Elabora{\c c}{\~a}o conjunta, refinamento e expans{\~a}o & Do estudante com aux{\'i}lio expl{\'i}cito de IA em trechos espec{\'i}ficos. \\
    \addlinespace
    \textbf{4} & IA Integral & Uso estrat{\'e}gico, abrangente e cont{\'i}nuo & Dirigido criticamente pelo estudante com forte presen{\c c}a de IA. \\
    \addlinespace
    \textbf{5} & Explora{\c c}{\~a}o de IA & Co-cria{\c c}{\~a}o avan{\c c}ada e inova{\c c}{\~a}o de fronteira & Parceria autoral entre estudante e IA em problemas complexos. \\
    \bottomrule
  \end{tabularx}
\end{table}

Duas propriedades estruturais da Escala AIAS s{\~a}o essenciais e devem orientar todo o desenho instrucional:
\begin{itemize}
  \item \textbf{Car{\'a}ter n{\~a}o hier{\'a}rquico}: Nenhum n{\'i}vel {\'e} intrinsecamente superior ou mais nobre do que outro. Uma atividade no N{\'i}vel~1 (Sem IA) pode demandar criatividade extrema e racioc{\'i}nio anal{\'i}tico profundo (por exemplo, um debate oral ou uma solu{\c c}{\~a}o manuscrita em sala), enquanto uma atividade no N{\'i}vel~4 pode focar em habilidades de auditoria cr{\'i}tica de c{\'o}digo ou revis{\~a}o de relat{\'o}rios. A escolha do n{\'i}vel deriva estritamente dos \textbf{objetivos de aprendizagem} visados;
  \item \textbf{Car{\'a}ter cumulativo}: N{\'i}veis superiores autorizam as pr{\'a}ticas dos n{\'i}veis inferiores, salvo proibi{\c c}{\~a}o expressa no enunciado. Por exemplo, uma atividade AIAS~3 admite que o aluno realize idea{\c c}{\~a}o inicial com a ferramenta (pr{\'a}tica do N{\'i}vel~2), avan{\c c}ando para o co-refinamento de se{\c c}{\~o}es espec{\'i}ficas.
\end{itemize}

\begin{atencaopedagogica}
\textbf{Declara{\c c}{\~a}o Pr{\'e}via Obrigat{\'o}ria}: O n{\'i}vel AIAS deve ser declarado de forma inequ{\'i}voca no enunciado da atividade e no plano de ensino \emph{antes} do in{\'i}cio da tarefa. Exigir do aluno uma postura em rela{\c c}{\~a}o {\`a} IA sem defini{\c c}{\~a}o formal pr{\'e}via compromete a seguran{\c c}a jur{\'i}dica e pedag{\'o}gica da avalia{\c c}{\~a}o.
\end{atencaopedagogica}

Ao planejar uma atividade, separe duas perguntas: a IA generativa est{\'a} \textbf{permitida} e a IA generativa est{\'a} \textbf{obrigat{\'o}ria}? O n{\'i}vel AIAS informa o papel da IA no processo, mas o enunciado deve esclarecer tamb{\'e}m quais evid{\^e}ncias ser{\~a}o avaliadas, como o uso ser{\'a} declarado e quais tecnologias assistivas ou adapta{\c c}{\~o}es institucionais permanecem autorizadas. ``Sem IA'' deve ser entendido como sem IA generativa na tarefa, n{\~a}o como proibi{\c c}{\~a}o de acessibilidade.
